\subsection{Cohomology and enveloping primitives in the completion}

Let $C=\bigoplus_{\beta\in Q_+}C_\beta$ be a charge-graded cochain
complex. The height product is taken in each cohomological degree.

\begin{proposition}[Cohomology of the split height completion]
\label{cohomology-completion}
For a charge-preserving differential,
\[
 H^p(\widehat C)=\prod_{\beta\in Q_+}H^p(C_\beta).
\]
A charge-preserving quasi-isomorphism therefore remains a
quasi-isomorphism after height completion.
\end{proposition}

\begin{proof}
Cycles and boundaries in a degreewise product are computed componentwise.
For boundaries, choose a preimage independently in every component.
Thus $Z^p(\widehat C)=\prod_\beta Z^p(C_\beta)$ and
$B^p(\widehat C)=\prod_\beta B^p(C_\beta)$.
The quotient of these products is the product of the quotients.
For a charge-preserving quasi-isomorphism of direct sums, each component
is a quasi-isomorphism.  Apply the displayed formula.
\end{proof}

The enveloping comparison has a particularly explicit completion.
Use the positive charge monoid $Q_+$ and its height as above, and assume
\[
 \mathfrak g=\bigoplus_{\beta\in Q_+\setminus\{0\}}\mathfrak g_\beta,
 \qquad d\mathfrak g_\beta\subseteq\mathfrak g_\beta,
 \qquad [\mathfrak g_\beta,\mathfrak g_\gamma]
 \subseteq\mathfrak g_{\beta+\gamma}.
\]
This charge is independent of cohomological degree.
Set $C=U(\mathfrak g)=\bigoplus_{\beta\in Q_+}C_\beta$, with $C_0=\kk$.
Its completion and completed tensor are those of
\eqref{completed-tensor}, taken degreewise. Multiplication is finite
convolution in each charge; coproduct, differential and antipode extend
chargewise by Theorem~\ref{tower-theorem}. The result is a topological
differential graded Hopf algebra in that specified tensor category.

\begin{proposition}[Primitives and cohomology in the charge completion]
\label{sp:completion}
With this charge grading, there are natural isomorphisms
\begin{align*}
 P(\widehat{U(\mathfrak g)})
   &=\prod_{\beta\ne0}\mathfrak g_\beta,\\
 H(\widehat{U(\mathfrak g)})
   &\cong\widehat{U(H(\mathfrak g))},\\
 H(P(\widehat{U(\mathfrak g)}))
   &\xrightarrow{\sim}P(H(\widehat{U(\mathfrak g)}))
    \cong\prod_{\beta\ne0}H(\mathfrak g_\beta).
\end{align*}
All products are degreewise.  The coproduct in cohomology uses the
same charge completion of the tensor product of the cohomology
components.  If $\mathfrak r$ is also charge-homogeneous, these
statements apply to the completion of
$U(\mathfrak g)/(\mathfrak r)\cong U(\mathfrak g/\mathfrak r)$.
\end{proposition}

\begin{proof}
The primitive equation holds precisely when it holds in each charge.
Theorem~\ref{sp:primitives} therefore proves the first identity.
Proposition~\ref{cohomology-completion} computes its cohomology
componentwise.
The isomorphism \eqref{sp:cohomology-map} preserves charge, as do
symmetrization and all the cycle cross products in its proof.
Taking these componentwise products gives the second isomorphism.
Applying Lemma~\ref{sp:kunnethtensor} to each pair of charge
components identifies the target of its cohomological coproduct.
Theorem~\ref{sp:cohomology}, charge by charge, now gives the last line.
The quotient assertion follows by applying the same argument to the
positive charge-graded Lie algebra $\mathfrak g/\mathfrak r$.
\end{proof}

For the abelian pair $du=v$, assign both generators charge one.
The completion is $\kk[[u]]\oplus \kk[[u]]v$, with $d f(u)=f'(u)v$.
Formal differentiation has kernel $\kk$ and is onto $\kk[[u]]v$:
the series $\sum_{m\geq0}a_mu^mv$ is the differential of
$\sum_{m\geq0}a_mu^{m+1}/(m+1)$.
Thus its cohomology is $\kk$, as the proposition asserts.
These are formal series in the stated height topology.
No analytic convergence is used.

The completion in Proposition~\ref{sp:completion} is
$\widehat{U(H(\mathfrak g))}$, formed after the algebraic enveloping
algebra.  It is not defined as $U(\prod_\beta H(\mathfrak g_\beta))$.
Its tensor product is also explicitly completed in charge.
The proposition makes no claim about a different topology or an
arbitrary completed tensor product.

